% titel: Blatt 1, A1: Messgenauigkeit (Mittelwert, Standardabweichung)
% skript: Mortimer Kap. 1 (chemie: Messwerte, Fehler)

% Teil I: Werkzeuge, kein Lösungsweg.

Mehrfachmessungen derselben Grösse: Mittelwert und Streuung bestimmen, um eine
Messreihe angemessen zu beschreiben.

\begin{werkzeugkasten}[Mittelwert und Standardabweichung]
Für $n$ Messwerte $x_1,\dots,x_n$:
\begin{equation}
  \bar x = \frac{1}{n}\sum_{i=1}^{n}x_i,\qquad
  s = \sqrt{\frac{1}{n-1}\sum_{i=1}^{n}(x_i-\bar x)^2}.
\end{equation}
$s$ ist die empirische (Stichproben-)Standardabweichung und misst, wie stark
die Werte im Schnitt um $\bar x$ streuen. Der Mittelwert ist die beste
Schätzung des wahren Werts, $s$ die zugehörige Unsicherheit.

\textit{Skript:} Mortimer Kap. 1 (Messung und Fehler), Übung 1/1.
\end{werkzeugkasten}

\begin{achtung}[Darauf achten]
\begin{itemize}
  \item Einheiten (hier g) überall mitschleppen; Ergebnis mit Einheit angeben.
  \item Zwischenwerte nicht runden, erst das Endergebnis auf sinnvolle
        Stellenzahl.
  \item $n-1$ im Nenner (Stichprobe), nicht $n$ (Grundgesamtheit) --- die
        Formel in der Vorlesung im Skript nachsehen.
  \item Standardabweichung rund um 0.01 g hier: deutet auf eine gute
        Waagen-Präzision bei ca. 26.5 g.
\end{itemize}
\end{achtung}
